% !TEX root = ../00_Main/Main.tex
\documentclass[../00_Main/Main.tex]{subfiles}
\begin{document}

\section{Comandos de reproducci\'on}\label{ap:repro}

Todos los comandos se ejecutan desde el directorio \texttt{h1/} del
repositorio \texttt{mPES}, con el entorno virtual \texttt{win\_mpes\_env}
activado, Python 3.12, las dependencias de
\texttt{utils/config/requirements.txt} y las variables
\texttt{VIRTUAL\_ENV}, \texttt{PYTHONIOENCODING=utf-8} y
\texttt{TF\_ENABLE\_ONEDNN\_OPTS=0} (Tabla~\ref{tab:repro-commands}). El
entrenamiento final lee los hiperpar\'ametros de
\texttt{inputs/best\_params.json}; la arquitectura del DQN recurrente y
del Transformer, y el n\'umero de episodios de todos los modelos, se
leen de \texttt{config/CONFIG.py}.

\begin{table}[!ht]
\centering
\caption[Comandos de reproducci\'on]{Comandos de optimizaci\'on bayesiana,
entrenamiento final y evaluaci\'on de cada paquete. \texttt{<pkg>} es
cualquiera de los cinco ensambles optimizables.}
\label{tab:repro-commands}
\small
\resizebox{\linewidth}{!}{%
\begin{tabular}{@{}llll@{}}
\toprule
\textbf{Paquete} & \textbf{Optimizaci\'on bayesiana} & \textbf{Entrenamiento final} & \textbf{Evaluaci\'on} \\
\midrule
\texttt{pes\_base} & --- & \texttt{python -m tabular.pes\_base.ext.train\_rl} & \texttt{python -m tabular.pes\_base} \\
\texttt{pes\_ql}   & \texttt{python -m tabular.pes\_ql.ext.optimize\_rl} & \texttt{python -m tabular.pes\_ql.ext.train\_rl} & \texttt{python -m tabular.pes\_ql} \\
\texttt{pes\_dql}  & \texttt{python -m tabular.pes\_dql.ext.optimize\_rl} & \texttt{python -m tabular.pes\_dql.ext.train\_rl} & \texttt{python -m tabular.pes\_dql} \\
\texttt{pes\_dqn}  & \texttt{python -m ml.pes\_dqn.ext.optimize\_dqn} & \texttt{python -m ml.pes\_dqn.ext.train\_dqn} & \texttt{python -m ml.pes\_dqn} \\
\texttt{pes\_rdqn} & \texttt{python -m ml.pes\_rdqn.ext.optimize\_rdqn} & \texttt{python -m ml.pes\_rdqn.ext.train\_rdqn} & \texttt{python -m ml.pes\_rdqn} \\
\texttt{pes\_a2c}  & \texttt{python -m ml.pes\_a2c.ext.optimize\_a2c} & \texttt{python -m ml.pes\_a2c.ext.train\_a2c} & \texttt{python -m ml.pes\_a2c} \\
\texttt{pes\_trf}  & \texttt{python -m ml.pes\_trf.ext.optimize\_tr} & \texttt{python -m ml.pes\_trf.ext.train\_transformer} & \texttt{python -m ml.pes\_trf} \\
\texttt{pes\_ens}  & --- & --- & \texttt{python -m ens.pes\_ens} \\
\texttt{<pkg>}     & \texttt{python -m ens.<pkg>.ext.optimize\_ens} & --- & \texttt{python -m ens.<pkg>} \\
\bottomrule
\end{tabular}
}%
\end{table}

\section{Ejecuci\'on de la evaluaci\'on}\label{ap:orchestrator}

La evaluaci\'on en los $22$ escenarios y las $5$ r\'eplicas \emph{held-out} se ejecuta desde \texttt{h1/}:

\begin{verbatim}
python -m general.scripts.benchmark run --suite both
python -m general.scripts.benchmark heldout
python -m general.scripts.analysis
python -m general.scripts.figures
python -m general.scripts.random_baseline
\end{verbatim}

El primer comando, reanudable, eval\'ua cada par (modelo, escenario) de
los dos grupos reemplazando temporalmente los archivos de entrada de
cada paquete por los del escenario. El segundo guarda en
\texttt{general/results/heldout/} una copia de las r\'eplicas y comprueba
que los $13$ modelos recibieron las mismas secuencias; cada r\'eplica
registra en \texttt{sampling\_distribution.json} las frecuencias de
sorteo, la semilla y las frecuencias efectivamente obtenidas. El tercero
calcula las matrices modelo $\times$ escenario y las comparaciones entre
pares; el cuarto genera las figuras, y el quinto, las del agente
aleatorio. Los resultados se guardan en
\texttt{general/results/<grupo>/}.

Las frecuencias de la Tabla~\ref{tab:ens-freq} se obtienen, desde la
ra\'iz del repositorio, con
\texttt{python writings/auxiliar/scripts/ensemble\_decisions.py}, que
repite la evaluaci\'on del ensamble ponderado, del consenso con
\emph{prior} y de la compuerta del Transformer, registra en cada decisi\'on
la acci\'on antes y despu\'es de cada regla y comprueba que las medias
coincidan con las de la evaluaci\'on principal. La
Tabla~\ref{tab:trf-vs-ens} se obtiene con
\texttt{python writings/auxiliar/scripts/trf\_vs\_ens.py}, que compara
cada ensamble con el Transformer a partir de los resultados guardados en
\texttt{general/results/}, y la Tabla~\ref{tab:heldout}, con
\texttt{python writings/auxiliar/scripts/heldout\_gap.py}. Los
an\'alisis de la Secci\'on~\ref{sec:res-posthoc} se obtienen, desde la
misma carpeta, con \texttt{weighted\_ens\_sensitivity.py} (sensibilidad),
\texttt{weighted\_ens\_oracle.py} (asignaci\'on \'optima y efecto de cada
regla) y \texttt{fixed\_rule.py} (regla fija en los $27$ escenarios).

\section{Mapas estad\'isticos complementarios}\label{ap:stat-maps}

Las figuras de este ap\'endice detallan, modelo por modelo, las
comparaciones estad\'isticas que el texto informa de forma agregada.

\begin{figure}[H]
  \centering
  \includegraphics[width=0.72\linewidth]{ind_13_pares_cohen_d.png}
  \caption[$d$ de Cohen entre pares (modelos individuales)]{Modelos
  individuales: $d$ de Cohen entre pares de modelos en los $21$ escenarios
  de generalizaci\'on.}
  \label{fig:pairwise-cohen}
\end{figure}

\begin{figure}[H]
  \centering
  \includegraphics[width=\linewidth]{ind_03_welch_logp_por_escenario.png}
  \caption[Welch por escenario (modelos individuales)]{Modelos
  individuales: $\log_{10} p$ de Welch de cada escenario frente a la
  referencia del mismo modelo; los valores menores que $-10$ se muestran
  como $\le -10$. En el panel inferior, las cinco columnas a la derecha de la l\'inea vertical son las r\'eplicas \emph{held-out}.}
  \label{fig:heatmap-welch}
\end{figure}

\begin{figure}[H]
  \centering
  \includegraphics[width=\linewidth]{ens_03_welch_logp_por_escenario.png}
  \caption[Welch por escenario (ensambles)]{Ensambles: $\log_{10} p$ de
  Welch de cada escenario frente a la referencia del mismo ensamble; los
  valores menores que $-10$ se muestran como $\le -10$. En el panel inferior, las cinco columnas a la derecha de la l\'inea vertical son las r\'eplicas \emph{held-out}.}
  \label{fig:ensemble-statistical-heatmaps}
\end{figure}

\begin{figure}[H]
  \centering
  \includegraphics[width=\linewidth]{ens_07_cohen_d_por_escenario.png}
  \caption[$d$ de Cohen por escenario (ensambles)]{Ensambles: $d$ de
  Cohen de cada escenario frente a la referencia del mismo ensamble
  (positivo: mejor que la referencia). En el panel inferior, las cinco columnas a la derecha de la l\'inea vertical son las r\'eplicas \emph{held-out}.}
  \label{fig:ensemble-cohen-scenario}
\end{figure}

\begin{figure}[!htbp]
  \centering
  \begin{subfigure}{\linewidth}
    \centering
    \includegraphics[width=0.95\linewidth]{ens_13_pares_cohen_d.png}
    \caption{$d$ de Cohen.}
  \end{subfigure}

  \begin{subfigure}{\linewidth}
    \centering
    \includegraphics[width=0.95\linewidth]{ens_12_pares_welch_logp.png}
    \caption{$\log_{10} p$ de Welch.}
  \end{subfigure}
  \caption[Comparaci\'on entre pares (ensambles)]{Ensambles:
  comparaci\'on entre pares en los $21$ escenarios de generalizaci\'on.}
  \label{fig:ensemble-pairwise}
\end{figure}

\section{Resultados por secuencia de cada modelo}\label{ap:per-model}

Las figuras de este ap\'endice muestran, para cada modelo individual, el
desempe\~no de cada secuencia de la referencia, su distribuci\'on y sus
estad\'isticos por bloque (Secci\'on~\ref{sec:res-individual}).

\begin{figure}[H]
  \centering
  \includegraphics[width=\linewidth]{PES_BASE_results.png}
  \caption[Resultados por secuencia: Q-Learning base]{Resultados por secuencia de Q-Learning base (\texttt{pes\_base}) en la referencia.}
  \label{fig:c-base}
\end{figure}

\begin{figure}[H]
  \centering
  \includegraphics[width=\linewidth]{PES_QL_results.png}
  \caption[Resultados por secuencia: Q-Learning]{Resultados por secuencia de Q-Learning (\texttt{pes\_ql}) en la referencia.}
  \label{fig:c-ql}
\end{figure}

\begin{figure}[H]
  \centering
  \includegraphics[width=\linewidth]{PES_DQL_results.png}
  \caption[Resultados por secuencia: Double Q-Learning]{Resultados por secuencia de Double Q-Learning (\texttt{pes\_dql}) en la referencia.}
  \label{fig:c-dql}
\end{figure}

\begin{figure}[H]
  \centering
  \includegraphics[width=\linewidth]{PES_DQN_results.png}
  \caption[Resultados por secuencia: DQN]{Resultados por secuencia de DQN (\texttt{pes\_dqn}) en la referencia.}
  \label{fig:c-dqn}
\end{figure}

\begin{figure}[H]
  \centering
  \includegraphics[width=\linewidth]{PES_RDQN_results.png}
  \caption[Resultados por secuencia: DQN recurrente]{Resultados por secuencia del DQN recurrente (\texttt{pes\_rdqn}) en la referencia.}
  \label{fig:c-rdqn}
\end{figure}

\begin{figure}[H]
  \centering
  \includegraphics[width=\linewidth]{PES_A2C_results.png}
  \caption[Resultados por secuencia: A2C]{Resultados por secuencia de A2C (\texttt{pes\_a2c}) en la referencia.}
  \label{fig:c-a2c}
\end{figure}

\begin{figure}[H]
  \centering
  \includegraphics[width=\linewidth]{PES_TRF_results.png}
  \caption[Resultados por secuencia: Transformer]{Resultados por secuencia del Transformer (\texttt{pes\_trf}) en la referencia.}
  \label{fig:c-trf}
\end{figure}

\biblio
\end{document}
